% Formation Engine -- XYZ File Formats (Edition 3)

% VSEPR-SIM 4.1.0  |  Compile: pdflatex (twice)

%

\documentclass[12pt, a4paper, oneside]{article}


\usepackage[top=1.0in, bottom=1.0in, left=1.15in, right=1.15in, headheight=14pt]{geometry}

\usepackage[utf8]{inputenc}

\usepackage[T1]{fontenc}

\usepackage{amsmath, amssymb, mathtools}

\usepackage{booktabs}

\usepackage{array}

\usepackage{tabularx}

\usepackage{enumitem}

\setlist{nosep, leftmargin=1.5em}

\usepackage{listings}

\usepackage{xcolor}

\usepackage[most]{tcolorbox}

\usepackage{mdframed}

\usepackage{multicol}


\lstset{basicstyle=\ttfamily\small, keywordstyle=\color{blue}, commentstyle=\color{gray},

  stringstyle=\color{red}, breaklines=true, columns=fullflexible, keepspaces=true,

  frame=single, backgroundcolor=\color{gray!8}, xleftmargin=0.5em, xrightmargin=0.5em,

  language=C++}


\usepackage{graphicx}

\usepackage{fancyhdr}

\pagestyle{fancy}

\fancyhf{}

\fancyhead[L]{\small\textit{Formation Engine Methodology}}

\fancyhead[R]{\small\textit{XYZ File Formats --- Edition 3}}

\fancyfoot[C]{\thepage}

\renewcommand{\headrulewidth}{0.4pt}


\usepackage{titlesec}

\titleformat{\section}{\Large\bfseries}{\S{}\thesection}{1em}{}

\titleformat{\subsection}{\large\bfseries}{\thesubsection}{1em}{}

\titleformat{\subsubsection}{\normalsize\bfseries}{\thesubsubsection}{1em}{}


\usepackage[colorlinks=true, linkcolor=black, citecolor=black, urlcolor=blue]{hyperref}


\newcommand{\file}[1]{\texttt{#1}}

\newcommand{\type}[1]{\texttt{#1}}

\newcommand{\field}[1]{\texttt{#1}}

\newcommand{\Ang}{\text{\AA}}

\newcommand{\kcalmol}{\text{kcal/mol}}


\title{%

  \textbf{Formation Engine Methodology}\\[6pt]

  \textbf{XYZ File Formats}\\[10pt]

  \large VSEPR-SIM Atomistic Engine\\

  Formats: \texttt{.xyz}\quad\texttt{.xyza}\quad\texttt{.xyzc}\quad\texttt{.xyzf}\quad\texttt{.xyzfull}\\[4pt]

  \normalsize Edition 3 --- Expanded Reference with Replay Format

}

\author{VSEPR-SIM Project}

\date{Version 4.1.0 --- \today}


\begin{document}

\maketitle


\begin{abstract}

All VSEPR-SIM molecular I/O uses plain-text, line-oriented files.

Four formats share a common structural grammar; each adds fields to

the preceding one.  This document is the expanded third-edition

reference for the \texttt{.xyz}, \texttt{.xyza}, \texttt{.xyzc},

\texttt{.xyzf}, and \texttt{.xyzfull} file formats, including grammar

specifications, field semantics, coordinate conventions, checkpoint and

trajectory usage, bounding-box encoding, implementation mapping, typical

observable ranges, and the rich replay container format.

\end{abstract}


\tableofcontents


% ============================================================

%  COMMENT BOX PAGES  (Edition 3 release notes + spec card)

% ============================================================

\newpage

\thispagestyle{fancy}


% ---- colour definitions ----

\definecolor{revblue}{RGB}{30,80,140}

\definecolor{revbluelight}{RGB}{220,232,245}

\definecolor{specgold}{RGB}{180,130,0}

\definecolor{specgoldlight}{RGB}{255,248,220}

\definecolor{warnred}{RGB}{160,30,30}

\definecolor{warnredlight}{RGB}{253,235,235}

\definecolor{okgreen}{RGB}{20,110,50}

\definecolor{okgreenlight}{RGB}{225,245,230}

\definecolor{coregray}{RGB}{60,60,70}

\definecolor{rulegray}{RGB}{180,185,195}


% ================================================================

\begin{tcolorbox}[

  enhanced, breakable=false,

  colback=revbluelight, colframe=revblue,

  boxrule=1.2pt, arc=4pt,

  title={\large\bfseries\color{white} What Is New in Edition 3 --- v4.1.0},

  attach boxed title to top left={yshift=-2mm, xshift=6mm},

  boxed title style={colback=revblue, colframe=revblue, arc=3pt},

  fonttitle=\bfseries, left=6pt, right=6pt, top=8pt, bottom=6pt

]


\textbf{Edition 3 adds one new format:} \texttt{.xyzfull} --- the rich replay and

render container.  All prior formats (\texttt{.xyz}, \texttt{.xyza}, \texttt{.xyzc},

\texttt{.xyzf}) are unchanged and fully backward-compatible.


\medskip

\begin{tcolorbox}[

  colback=white, colframe=rulegray,

  boxrule=0.6pt, arc=2pt, left=4pt, right=4pt, top=4pt, bottom=4pt

]

\textbf{New in this edition:}

\begin{itemize}

  \item \textbf{\S{}6 --- Full Replay XYZ (\texttt{.xyzfull}):}  complete grammar,

        data model, energy layer convention, STEP proxy tag, orientation vector

        spec, streaming writer API, shell commands, and validation rules.

  \item \textbf{Format hierarchy table} updated to include \texttt{.xyzfull}.

  \item \textbf{Feature comparison matrix} extended to five formats (15 feature rows).

  \item \textbf{Implementation mapping table} updated with replay renderer files.

  \item \textbf{Reader / Writer API patterns} updated with \texttt{XYZFullReader}

        and \texttt{XYZFullStreamWriter}.

  \item \textbf{\S{}A --- Optional \texttt{INFO} Block and Compute Kernel Tags:}

        \texttt{INFO} extended as an optional tail block for all formats; kernel

        utilisation metadata defined.

\end{itemize}

\end{tcolorbox}


\medskip

\begin{tcolorbox}[

  colback=warnredlight, colframe=warnred,

  boxrule=0.6pt, arc=2pt, left=4pt, right=4pt, top=4pt, bottom=4pt

]

\textbf{\color{warnred}Critical distinction to memorise:}

\texttt{.xyzfull} is a \emph{replay artifact}, not a checkpoint.

Restarting physics from render crumbs is not supported.

\end{tcolorbox}


\end{tcolorbox}


\bigskip


% ================================================================

% PAGE 2 -- .xyzfull SPEC CARD

% ================================================================

\newpage

\thispagestyle{fancy}


\begin{tcolorbox}[

  enhanced,

  colback=gray!5, colframe=coregray,

  boxrule=1.2pt, arc=4pt,

  title={\large\bfseries\color{white} \texttt{.xyzfull} --- At-a-Glance Spec Card},

  attach boxed title to top left={yshift=-2mm, xshift=6mm},

  boxed title style={colback=coregray, colframe=coregray, arc=3pt},

  left=6pt, right=6pt, top=8pt, bottom=6pt

]


\begin{multicols}{2}


% --- LEFT COLUMN ---

\begin{tcolorbox}[

  colback=revbluelight, colframe=revblue,

  boxrule=0.7pt, arc=3pt, left=4pt, right=4pt, top=3pt, bottom=3pt,

  title={\bfseries\small File Structure},

  fonttitle=\small\bfseries, attach title to upper, after title={\\[2pt]}

]

\small

\texttt{XYZFULL}\\

\texttt{version 1}\\

\texttt{units angstrom fs kcal/mol}\\

\texttt{mode <string>}\\

\texttt{energy\_layers <N> <names...>}\\

\texttt{static\_count <S>}\\

\texttt{dynamic\_count <D>}\\

\texttt{properties ...}\\[4pt]

\texttt{STATIC}\\

\texttt{<id> <type> <sym> x y z q cg\_id}\\[4pt]

\texttt{FRAMES}\\

\texttt{FRAME <k> time <t>}\\

\texttt{<dynamic rows>}\\

\texttt{END\_FRAME}\\[4pt]

\texttt{INFO}\\

\texttt{<free text>}\\

\texttt{TOTAL\_FRAMES <N>}\\

\texttt{DT <dt>}

\end{tcolorbox}


\medskip


\begin{tcolorbox}[

  colback=okgreenlight, colframe=okgreen,

  boxrule=0.7pt, arc=3pt, left=4pt, right=4pt, top=3pt, bottom=3pt,

  title={\bfseries\small Energy Layers},

  fonttitle=\small\bfseries, attach title to upper, after title={\\[2pt]}

]

\small

\begin{tabular}{ll}

\texttt{e0} & \texttt{U\_total} \\

\texttt{e1} & \texttt{U\_bond} \\

\texttt{e2} & \texttt{U\_vdw} \\

\texttt{e3} & \texttt{U\_coul} \\

\texttt{e4} & \texttt{U\_ext} \\

\texttt{e5} & \texttt{U\_thermal} \\

\end{tabular}\\[4pt]

Count: $k \in \{4, 5, 6\}$.  Declare in header.\\

Do not invent arbitrary counts.

\end{tcolorbox}


\columnbreak


% --- RIGHT COLUMN ---

\begin{tcolorbox}[

  colback=specgoldlight, colframe=specgold,

  boxrule=0.7pt, arc=3pt, left=4pt, right=4pt, top=3pt, bottom=3pt,

  title={\bfseries\small Key Constraints},

  fonttitle=\small\bfseries, attach title to upper, after title={\\[2pt]}

]

\small

\begin{itemize}[leftmargin=1em]

  \item File \emph{must} start with \texttt{XYZFULL}

  \item Last two lines: \texttt{TOTAL\_FRAMES} then \texttt{DT}

  \item Parsed frame count $=$ \texttt{TOTAL\_FRAMES}

  \item Dynamic IDs unique within frame

  \item Static IDs $\neq$ dynamic IDs

  \item $\|\mathbf{o}\| > 10^{-12}$ (orientation vector)

  \item Energy layer count: 4, 5, or 6 only

  \item Unknown \texttt{step\_tag}: warn, not fatal

\end{itemize}

\end{tcolorbox}


\medskip


\begin{tcolorbox}[

  colback=white, colframe=rulegray,

  boxrule=0.7pt, arc=3pt, left=4pt, right=4pt, top=3pt, bottom=3pt,

  title={\bfseries\small CG Proxy Shapes},

  fonttitle=\small\bfseries, attach title to upper, after title={\\[2pt]}

]

\small

\begin{tabular}{ll}

\texttt{bead}    & sphere / dot \\

\texttt{rod}     & capsule / line \\

\texttt{plate}   & rectangle proxy \\

\texttt{gear}    & low-poly wheel \\

\texttt{pipe}    & tube segment \\

\texttt{crystal} & lattice proxy \\

\end{tabular}\\[4pt]

Not CAD.  Visual truth only.

\end{tcolorbox}


\medskip


\begin{tcolorbox}[

  colback=warnredlight, colframe=warnred,

  boxrule=0.7pt, arc=3pt, left=4pt, right=4pt, top=3pt, bottom=3pt,

  title={\bfseries\small NOT a checkpoint},

  fonttitle=\small\bfseries, attach title to upper, after title={\\[2pt]}

]

\small

\texttt{.xyzfull} $\neq$ restartable state.\\

Use \texttt{.xyzc} to restart MD.\\

Use \texttt{XYZFullStreamWriter} for\\

long runs --- not RAM.

\end{tcolorbox}


\end{multicols}


\end{tcolorbox}


\newpage


% ============================================================

%  END COMMENT BOX PAGES

% ============================================================


\newpage


% ================================================

\section{Overview}

% ================================================


The VSEPR-SIM atomistic engine reads and writes molecular geometry,

dynamics, and checkpoint data using four progressively richer

plain-text formats.  Each format extends the previous one by appending

per-atom or per-frame metadata while preserving backward compatibility

with simple line-oriented parsers.


\begin{table}[htbp]

  \centering

  \caption{Format hierarchy.}

  \label{tab:hierarchy}

  \begin{tabular}{llp{8cm}}

    \toprule

    Extension & Name & Use case \\

    \midrule

    \texttt{.xyz}     & Standard XYZ     & Static geometry snapshot \\

    \texttt{.xyza}    & Extended XYZ     & Charges, velocities, forces \\

    \texttt{.xyzc}    & Checkpoint XYZ   & Restartable MD state \\

    \texttt{.xyzf}    & Formation XYZ    & Multi-frame trajectory \\

    \texttt{.xyzfull} & Full Replay XYZ  & Rich replay/render container \\

    \bottomrule

  \end{tabular}

\end{table}


\subsection{Unit Conventions}


All formats share a single, non-negotiable unit system:


\begin{table}[htbp]

  \centering

  \caption{Unit system (all formats).}

  \label{tab:units}

  \begin{tabular}{lll}

    \toprule

    Quantity & Unit & SI equivalent \\

    \midrule

    Coordinates   & \AA{} (angstrom) & $10^{-10}$\,m \\

    Energy        & kcal/mol         & $4.184 \times 10^{3}$\,J/mol \\

    Velocity      & \AA/fs           & $10^{5}$\,m/s \\

    Force         & kcal/(mol\,\AA)  & $6.947 \times 10^{13}$\,N/mol \\

    Charge        & elementary $e$   & $1.602 \times 10^{-19}$\,C \\

    Time          & fs               & $10^{-15}$\,s \\

    Temperature   & K                & kelvin \\

    \bottomrule

  \end{tabular}

\end{table}


\subsection{Coordinate System}


Right-handed Cartesian $(x, y, z)$ in \AA.  Origin is arbitrary;

the conventional choice is near the centre of mass.  Coordinate

precision: six decimal places ($\Rightarrow$ resolution $\approx 0.1$\,pm).


% ================================================

\section{Standard XYZ (\texttt{.xyz})}

\label{sec:xyz}

% ================================================


\subsection{Grammar}


\begin{lstlisting}[language={},frame=single]

<N>

<comment>

<sym> <x> <y> <z>

...  (N lines)

\end{lstlisting}


\begin{itemize}

  \item \textbf{Line 1}: atom count $N$ (integer $\ge 1$).

  \item \textbf{Line 2}: free-text comment (may be empty).

  \item \textbf{Lines 3--$(N{+}2)$}: element symbol + three Cartesian

        coordinates, whitespace-separated.

\end{itemize}


\subsection{Comment Line Convention}


The comment line follows a structured-but-optional pattern:


\begin{lstlisting}[language={},frame=single]

<name> | E = <val> kcal/mol | <notes>

\end{lstlisting}


Fields are pipe-delimited.  The energy key \texttt{E =} is recognised

by the reader for energy-tracking contexts.  Additional free-text

notes are permitted after the second pipe.


\subsection{Example: Water}


\begin{lstlisting}[language={},frame=single]

3

H2O | E = -6.245 kcal/mol

O   0.000000   0.000000   0.119262

H   0.000000   0.763239  -0.477049

H   0.000000  -0.763239  -0.477049

\end{lstlisting}


\subsection{Parsing Notes}


\begin{itemize}

  \item Element symbols are case-sensitive: \texttt{He} is helium,

        \texttt{HE} is rejected.

  \item Trailing whitespace on atom lines is ignored.

  \item Blank lines between the comment line and the first atom line

        are \emph{not} permitted; the reader expects $N$ consecutive

        atom lines starting at line 3.

  \item Coordinate fields must parse as IEEE~754 double-precision

        floating-point values.

\end{itemize}


% ================================================

\section{Extended XYZ (\texttt{.xyza})}

\label{sec:xyza}

% ================================================


Adds per-atom scalar and vector properties after the $(x, y, z)$ fields.


\subsection{Grammar}


\begin{lstlisting}[language={},frame=single]

<N>

<comment> properties="<list>","<sym> <x> <y> <z> [<q> <vx> <vy> <vz> <fx> <fy> <fz> <e>]

\end{lstlisting}


The \texttt{properties} key declares which optional columns are present.

Column order is fixed; omission is permitted only from the right.


\subsection{Property Columns}


\begin{table}[htbp]

  \centering

  \caption{Extended property columns (fixed order).}

  \label{tab:xyza_cols}

  \begin{tabular}{llll}

    \toprule

    Key & Symbol & Unit & Columns \\

    \midrule

    \texttt{charge}   & $q_i$          & $e$                   & 1 \\

    \texttt{velocity} & $\mathbf{v}_i$ & \AA/fs                & 3 \\

    \texttt{force}    & $\mathbf{F}_i$ & kcal/(mol\,\AA)       & 3 \\

    \texttt{energy}   & $\varepsilon_i$& kcal/mol              & 1 \\

    \bottomrule

  \end{tabular}

\end{table}


\subsection{Properties Declaration Syntax}


The \texttt{properties} value is a colon-separated list of keys:


\begin{lstlisting}[language={},frame=single]

properties="charge:velocity"

properties="charge:velocity:force:energy"

\end{lstlisting}


Unrecognised keys trigger a parse warning but do not abort the read.


\subsection{Example: H\texorpdfstring{$_2$}{2}O with Charges and Velocities}


\begin{lstlisting}[language={},frame=single]

3

Water MD frame 1 properties="charge:velocity"

O   0.00   0.00   0.12  -0.834   0.00   0.00   0.00

H   0.00   0.76  -0.48   0.417   0.12   0.08  -0.03

H   0.00  -0.76  -0.48   0.417  -0.12   0.08  -0.03

\end{lstlisting}


\subsection{Charge Semantics}


The charge field $q_i$ represents the effective partial charge used in

Coulomb interactions.  It is \emph{not} an oxidation state.  For

Drude-polarisable models, $q_i$ is the core charge; the Drude

particle charge is stored separately in the polarisation subsystem.


% ================================================

\section{Checkpoint XYZ (\texttt{.xyzc})}

\label{sec:xyzc}

% ================================================


Extends \texttt{.xyza} with a simulation-state header block required

to restart an interrupted MD run without loss of trajectory continuity.


\subsection{Header Block}


The header block appears \emph{before} the atom-count line, delimited

by \texttt{CHECKPOINT} / \texttt{END\_CHECKPOINT} markers:


\begin{lstlisting}[language={},frame=single]

CHECKPOINT

step       <integer>

time       <float> fs

dt         <float> fs

T_target   <float> K

seed       <integer>

box        <ax> <ay> <az>

END_CHECKPOINT

<N>

<comment>

<atom lines with charge + velocity + force>

\end{lstlisting}


\subsection{Header Field Semantics}


\begin{table}[htbp]

  \centering

  \caption{Checkpoint header fields.}

  \label{tab:xyzc_header}

  \begin{tabular}{llp{7.5cm}}

    \toprule

    Field & Type & Meaning \\

    \midrule

    \texttt{step}      & integer & MD step number at checkpoint \\

    \texttt{time}      & float   & Accumulated simulation time (fs) \\

    \texttt{dt}        & float   & Integrator timestep (fs) \\

    \texttt{T\_target} & float   & Thermostat set-point (K); 0 = NVE ensemble \\

    \texttt{seed}      & integer & RNG seed for reproducible continuation \\

    \texttt{box}       & 3 floats & Simulation box side lengths (\AA);

                                    omit line entirely for open boundary \\

    \bottomrule

  \end{tabular}

\end{table}


\subsection{Restart Protocol}


When the engine loads a \texttt{.xyzc} file:


\begin{enumerate}

  \item The header is parsed; integrator state restored (step, time, dt,

        thermostat target, RNG state).

  \item The atom block is read as a standard \texttt{.xyza} frame.

  \item Neighbour lists and bond tables are reconstructed from coordinates;

        they are \emph{never} serialised into the checkpoint.

  \item If a \texttt{box} line is present, PBC is activated.

\end{enumerate}


\subsection{Checkpoint vs.\ Trajectory}


A \texttt{.xyzc} file contains exactly one frame.  For multi-frame

output, use \texttt{.xyzf} (\S\ref{sec:xyzf}).


% ================================================

\section{Formation XYZ (\texttt{.xyzf})}

\label{sec:xyzf}

% ================================================


A multi-frame format: one or more \texttt{.xyz} or \texttt{.xyza} blocks

concatenated in a single file.  Used for minimisation traces, NVT

trajectories, and reaction paths.


\subsection{Grammar}


\begin{lstlisting}[language={},frame=single]

<N>           <- frame 1

<comment>

<atom lines>

<N>           <- frame 2

<comment>

<atom lines>

...

\end{lstlisting}


Each frame is self-contained and independently parseable.


\subsection{Comment Line Convention (Multi-Frame)}


The comment line of frame $k$ typically records the step index and

instantaneous energy:


\begin{lstlisting}[language={},frame=single]

Step 0 | E = -12.30 kcal/mol | FIRE iter

Step 1 | E = -13.11 kcal/mol | FIRE iter

\end{lstlisting}


\subsection{Frame Independence}


Each frame must be independently parseable:


\begin{itemize}

  \item Atom count $N$ may vary between frames.

  \item The \texttt{properties} declaration may differ between frames.

  \item Readers must re-parse the header of each frame; caching the

        first frame's column layout globally is a bug.

\end{itemize}


\subsection{Typical Uses}


\begin{table}[htbp]

  \centering

  \caption{Common \texttt{.xyzf} applications.}

  \label{tab:xyzf_uses}

  \begin{tabular}{lp{8cm}}

    \toprule

    Application & Comment line pattern \\

    \midrule

    FIRE minimisation trace  & \texttt{Step <k> | E = <val> kcal/mol | FIRE iter} \\

    NVT trajectory dump      & \texttt{Step <k> | E = <val> kcal/mol | T = <T> K} \\

    Reaction path (NEB)      & \texttt{Image <k> | E = <val> kcal/mol | dE = <dE>} \\

    Conformer scan           & \texttt{Conf <k> | E = <val> kcal/mol | <dihedral>} \\

    \bottomrule

  \end{tabular}

\end{table}


% ================================================

\section{Bounding Box and Periodic Boundary Conditions}

\label{sec:pbc}

% ================================================


When PBC are active, box geometry is stored in the \texttt{.xyzc} header

(\S\ref{sec:xyzc}) or as comment-line keys in any format.


\subsection{Comment-Line Encoding}


\begin{lstlisting}[language={},frame=single]

Lattice="ax 0 0  0 ay 0  0 0 az" pbc="T T T"

\end{lstlisting}


\begin{itemize}

  \item Box vectors in \AA, flattened $3\times 3$ row-major matrix.

        Orthorhombic: diagonal entries only; triclinic: full matrix.

  \item \texttt{pbc} flags: \texttt{T} = periodic, \texttt{F} = open.

  \item Compatible with ASE Extended XYZ convention.

\end{itemize}


\subsection{Box Vector Interpretation}


For an orthorhombic cell with side lengths $a_x, a_y, a_z$:

\begin{equation}

  \mathbf{L} =

  \begin{pmatrix}

    a_x & 0   & 0   \\

    0   & a_y & 0   \\

    0   & 0   & a_z

  \end{pmatrix}

  \label{eq:lattice}

\end{equation}

The cell volume is $V = \det(\mathbf{L}) = a_x \cdot a_y \cdot a_z$.


\subsection{Minimum Image Convention}


Under PBC, the engine applies the minimum-image convention:

\begin{equation}

  \Delta r_\alpha = r_\alpha^{(j)} - r_\alpha^{(i)}

    - L_\alpha \cdot \mathrm{round}\!\left(

      \frac{r_\alpha^{(j)} - r_\alpha^{(i)}}{L_\alpha}

    \right),

  \qquad \alpha \in \{x, y, z\}

  \label{eq:mic}

\end{equation}

where $L_\alpha$ is the box length along axis $\alpha$.


% ================================================

\section{Bond Auto-Detection}

\label{sec:bonds}

% ================================================


Bond connectivity is inferred from geometry, never stored in the file.

The \texttt{detect\_bonds} function uses per-element covalent radii:

\begin{equation}

  d_{ij} < s \cdot \bigl(r_{\text{cov}}(Z_i) + r_{\text{cov}}(Z_j)\bigr)

  \quad \Longrightarrow \quad \text{bond}(i, j)

  \label{eq:bond}

\end{equation}

where $s$ is the scale factor (default $s = 1.2$).


\subsection{Design Rationale}


Bonds are derived data, not ground truth.  Stale bond tables are a

silent-corruption vector.  The \S{}0 container philosophy formalises

this: caches never become truth.


% ================================================

\section{Implementation Mapping}

\label{sec:impl}

% ================================================


\begin{table}[htbp]

  \centering

  \caption{Source files implementing the XYZ format family.}

  \label{tab:impl}

  \begin{tabular}{lp{8.5cm}}

    \toprule

    File & Role \\

    \midrule

    \file{include/io/xyz\_format.hpp}

      & Public API: \type{XYZReader}, \type{XYZWriter},

        \type{XYZAReader}, \type{XYZAWriter} \\

    \file{src/io/xyz\_format.cpp}

      & Reader / writer implementation \\

    \file{include/thermal/xyzc\_format.hpp}

      & Checkpoint reader/writer declarations \\

    \file{src/thermal/xyzc\_format.cpp}

      & Checkpoint reader/writer implementation \\

    \file{include/io/xyzfull\_format.hpp}

      & Replay container: \type{XYZFullReader}, \type{XYZFullWriter},

        \type{XYZFullStreamWriter} \\

    \file{src/io/xyzfull\_format.cpp}

      & Full replay reader / writer / streaming writer \\

    \file{include/render/xyzfull\_replay.hpp}

      & Low-poly replay renderer declarations \\

    \file{src/render/xyzfull\_replay.cpp}

      & Replay renderer and CG proxy reconstruction \\

    \bottomrule

  \end{tabular}

\end{table}


All reader and writer types reside in namespace \type{vsepr::io}.


\subsection{Reader API Pattern}


\begin{lstlisting}

// Standard XYZ

auto frame = vsepr::io::XYZReader::read("molecule.xyz");

// Extended XYZ

auto frame = vsepr::io::XYZAReader::read("dynamics.xyza");

// Checkpoint

auto state = vsepr::io::XYZCReader::read("restart.xyzc");

// Formation (multi-frame)

auto traj  = vsepr::io::XYZFReader::read_all("trajectory.xyzf");

// or frame-by-frame:

vsepr::io::XYZFReader reader("trajectory.xyzf");

while (auto frame = reader.next()) { /* ... */ }

// Full replay container

auto replay = vsepr::io::XYZFullReader::read("replay.xyzfull");

\end{lstlisting}


\subsection{Writer API Pattern}


\begin{lstlisting}

vsepr::io::XYZWriter::write("output.xyz", atoms, comment);

vsepr::io::XYZAWriter::write("output.xyza", atoms, props, comment);

vsepr::io::XYZCWriter::write("checkpoint.xyzc", state);

vsepr::io::XYZFWriter writer("trajectory.xyzf", /*append=*/true);

writer.write_frame(atoms, comment);

// Full replay: use streaming writer for long simulations

vsepr::io::XYZFullStreamWriter sw("replay.xyzfull");

sw.write_header(header);

sw.write_static_block(statics);

sw.begin_frames();

for (const auto& frame : frames) sw.write_frame(frame);

sw.write_footer(info, total_frames, dt);

\end{lstlisting}


% ================================================

\section{Typical Energy and Observable Ranges}

\label{sec:ranges}

% ================================================


\begin{table}[htbp]

  \centering

  \caption{Typical observable magnitudes for sanity checking.}

  \label{tab:ranges}

  \begin{tabular}{llr}

    \toprule

    Observable & Typical range & Unit \\

    \midrule

    Total bonded energy          & $-10^2$ to $-10^5$   & kcal/mol \\

    Relative conformer energy    & 0.1 -- 10            & kcal/mol \\

    Activation barrier           & 5 -- 50              & kcal/mol \\

    Force (unrelaxed atom)       & 1 -- 100             & kcal/(mol\,\AA) \\

    Force convergence criterion  & $< 0.01$             & kcal/(mol\,\AA) \\

    Velocity (300\,K, C atom)    & $\sim 0.02$          & \AA/fs \\

    Partial charge (O in water)  & $-0.82$ to $-0.84$   & $e$ \\

    Partial charge (H in water)  & $+0.41$ to $+0.42$   & $e$ \\

    \bottomrule

  \end{tabular}

\end{table}


These ranges serve as first-pass sanity checks.  Values outside these

bounds in production output should trigger audit warnings.


% ================================================

\section{Integration with the Identity--State Framework}

\label{sec:s0}

% ================================================


The XYZ formats serve as the serialisation layer for the particle

identity vector $\mathbf{I}_i$ defined in \S{}0 of the Formation

Engine methodology.


\begin{table}[htbp]

  \centering

  \caption{Identity vector components and their file-format encoding.}

  \label{tab:s0_map}

  \begin{tabular}{llp{6.5cm}}

    \toprule

    Component & Symbol & File encoding \\

    \midrule

    Nuclear identity     & $Z_i$       & Element symbol (looked up to $Z$) \\

    Mass identity        & $A_i$       & Implicit via $Z_i$ \\

    Charge participation & $Q_i$       & \texttt{charge} column in \texttt{.xyza} \\

    Structural role      & $\Sigma_i$  & Not serialised; inferred at load \\

    Stability class      & $\Lambda_i$ & Not serialised; computed post-load \\

    Provenance memory    & $\Theta_i$  & Comment-line metadata \\

    \bottomrule

  \end{tabular}

\end{table}


$\mathcal{K}_{\text{geom}}$ maps to the lattice matrix from the

$\texttt{Lattice}$ key or \texttt{box} header.

$\mathcal{K}_{\text{comp}}$ is reconstructed at load from element

composition and geometry.  The source-layer flag $F_K \in \{Z, A, C\}$

maps \texttt{Z}$\to$\texttt{.xyz}, \texttt{A}$\to$\texttt{.xyza},

\texttt{C}$\to$\texttt{.xyzc}.


% ================================================

\section{Format Comparison Matrix}

\label{sec:matrix}

% ================================================


\begin{table}[htbp]

  \centering

  \caption{Feature comparison across formats.}

  \label{tab:matrix}

  \begin{tabular}{lccccc}

    \toprule

    Feature & \texttt{.xyz} & \texttt{.xyza} & \texttt{.xyzc} & \texttt{.xyzf} & \texttt{.xyzfull} \\

    \midrule

    Atom count + coordinates       & \checkmark & \checkmark & \checkmark & \checkmark & \checkmark \\

    Partial charges                &            & \checkmark & \checkmark & opt.       & static     \\

    Velocities                     &            & \checkmark & \checkmark & opt.       & \checkmark \\

    Forces                         &            & \checkmark & \checkmark & opt.       &            \\

    Per-atom energy                &            & \checkmark & \checkmark & opt.       & layered    \\

    Integrator state               &            &            & \checkmark &            &            \\

    Thermostat / RNG               &            &            & \checkmark &            &            \\

    Box geometry (header)          &            &            & \checkmark &            &            \\

    Box geometry (comment key)     & \checkmark & \checkmark & \checkmark & \checkmark &            \\

    Multi-frame                    &            &            &            & \checkmark & \checkmark \\

    Static / dynamic separation    &            &            &            &            & \checkmark \\

    CG orientation vector          &            &            &            &            & \checkmark \\

    Energy layer decomposition     &            &            &            &            & \checkmark \\

    STEP proxy tag                 &            &            &            &            & \checkmark \\

    Low-poly render support        &            &            &            &            & \checkmark \\

    \bottomrule

  \end{tabular}

\end{table}


% ================================================

\section{Error Handling and Validation}

\label{sec:errors}

% ================================================


Readers enforce these validation rules at parse time:


\begin{enumerate}

  \item \textbf{Atom count mismatch.}  Fewer than $N$ atom lines

        $\Rightarrow$ \type{XYZParseError::INCOMPLETE\_FRAME}.

  \item \textbf{Unknown element.}  Unrecognised symbol

        $\Rightarrow$ \type{XYZParseError::UNKNOWN\_ELEMENT}.

        Two-letter symbols must match the periodic table (case-sensitive).

  \item \textbf{Malformed coordinates.}  Non-numeric field

        $\Rightarrow$ \type{XYZParseError::BAD\_COORDINATE}.

  \item \textbf{Column count mismatch.}  Declared \texttt{properties}

        implies $k$ columns but line has fewer: warn and zero-fill

        (fail-soft for velocities/forces; fail-hard for coordinates).

  \item \textbf{Checkpoint integrity.}  Missing

        \texttt{CHECKPOINT}/\texttt{END\_CHECKPOINT} block

        $\Rightarrow$ \type{XYZCParseError::NO\_HEADER}.

\end{enumerate}


% ================================================

\section{Full Replay XYZ (\texttt{.xyzfull})}

\label{sec:xyzfull}

% ================================================


The \texttt{.xyzfull} format is the terminal form of the XYZ family: a rich replay

and render container that captures bead/CG dynamics, orientation state, energy

layers, and STEP proxy tags in a single streamable file.


\textbf{It is not a checkpoint.  It is not a canonical simulation state.

It is a replay artifact.}  Restarting physics from \texttt{.xyzfull} data is

not supported and would be unsound.


\subsection{Design Rationale}


\begin{itemize}

  \item \texttt{.xyzfull} separates \emph{static} particles (walls, fixtures)

        from \emph{dynamic} particles (beads, CG bodies).  Statics are stored

        once; dynamic state is stored per frame.

  \item Orientation vectors $(\mathbf{o})$ enable low-poly proxy reconstruction

        of CG shapes without embedding CAD geometry per frame.

  \item Energy layers are declared in the header and stored inline, supporting

        4--6 named components.

  \item STEP proxy tags link each bead to a part identifier for

        renderer-side shape look-up.

\end{itemize}


\subsection{Grammar}


\begin{lstlisting}[language={},frame=single]

XYZFULL

version 1

units angstrom fs kcal/mol

mode <mode_string>

energy_layers <N> <name0> ... <nameN-1>

static_count <S>

dynamic_count <D>

properties id type sym x y z vx vy vz ox oy oz cg_id step_tag e0 [e1 ...]


STATIC

<id> <type> <sym> <x> <y> <z> <q> <cg_id>

...  (S lines)


FRAMES

FRAME <k> time <t>

<id> <type> <sym> <x> <y> <z> <vx> <vy> <vz> <ox> <oy> <oz> <cg_id> <step_tag> <e0> [...]

...  (D lines)

END_FRAME

...


INFO

<free text>

TOTAL_FRAMES <N>

DT <dt>

\end{lstlisting}


\subsection{File-Level Rules}


\begin{itemize}

  \item File must begin with the token \texttt{XYZFULL} on the first line.

  \item \texttt{version} must be a supported integer (currently 1).

  \item \texttt{energy\_layers} declares the count ($k \in \{4, 5, 6\}$) followed

        by that many whitespace-separated layer names.

  \item The final two lines of the file must be

        \texttt{TOTAL\_FRAMES \textit{N}} and \texttt{DT \textit{dt}}.

  \item Parsed frame count must equal \texttt{TOTAL\_FRAMES}.

\end{itemize}


\subsection{Energy Layer Convention}


\begin{table}[htbp]

  \centering

  \caption{Recommended 6-layer energy decomposition for \texttt{.xyzfull}.}

  \label{tab:xyzfull_energy}

  \begin{tabular}{lll}

    \toprule

    Index & Key & Description \\

    \midrule

    \texttt{e0} & \texttt{U\_total}   & Total potential energy \\

    \texttt{e1} & \texttt{U\_bond}    & Bonded interactions \\

    \texttt{e2} & \texttt{U\_vdw}     & van der Waals \\

    \texttt{e3} & \texttt{U\_coul}    & Coulomb/electrostatic \\

    \texttt{e4} & \texttt{U\_ext}     & External field \\

    \texttt{e5} & \texttt{U\_thermal} & Thermal / thermostat contribution \\

    \bottomrule

  \end{tabular}

\end{table}


A reduced 4-layer mode (\texttt{e0}=\texttt{U\_total}, \texttt{e1}=\texttt{U\_internal},

\texttt{e2}=\texttt{U\_external}, \texttt{e3}=\texttt{U\_state}) is also valid.

Declare the chosen scheme in the header:


\begin{lstlisting}[language={},frame=single]

energy_layers 6 U_total U_bond U_vdw U_coul U_ext U_thermal

\end{lstlisting}


Do not create arbitrary layer counts.  The instinct toward infinite

extensibility is natural and will produce vapor you then have to debug.


\subsection{Orientation Vector and STEP Proxy}


The orientation vector $(o_x, o_y, o_z)$ must be unit-normalizable:

\begin{equation}

  \|\mathbf{o}\| = \sqrt{o_x^2 + o_y^2 + o_z^2} > 10^{-12}

  \label{eq:orient}

\end{equation}


Violation produces a warning; the frame is not discarded.


The \texttt{step\_tag} + \texttt{cg\_id} pair is the lookup key for renderer-side

proxy reconstruction:


\begin{table}[htbp]

  \centering

  \caption{CG proxy shape map (low-poly, not full CAD).}

  \label{tab:xyzfull_proxy}

  \begin{tabular}{ll}

    \toprule

    Bead type & Proxy geometry \\

    \midrule

    \texttt{bead}    & Sphere / dot \\

    \texttt{rod}     & Capsule / line \\

    \texttt{plate}   & Rectangle proxy \\

    \texttt{gear}    & Low-poly wheel \\

    \texttt{pipe}    & Tube segment \\

    \texttt{crystal} & Lattice proxy \\

    \bottomrule

  \end{tabular}

\end{table}


Unknown \texttt{step\_tag} values are permitted but generate a parse warning.

Storing full CAD geometry per frame is explicitly forbidden by this specification,

as the universe catches fire and the file grows without bound.


\subsection{Minimal Example}


\begin{lstlisting}[language={},frame=single]

XYZFULL

version 1

units angstrom fs kcal/mol

mode bead_cg_replay

energy_layers 6 U_total U_bond U_vdw U_coul U_ext U_thermal

static_count 2

dynamic_count 2

properties id type sym x y z vx vy vz ox oy oz cg_id step_tag e0 e1 e2 e3 e4 e5


STATIC

900 wall X 0.000000 0.000000 0.000000 0.0 -1

901 wall X 10.00000 0.000000 0.000000 0.0 -1


FRAMES

FRAME 0 time 0.000000

1 bead C 0.000000 0.000000 0.000000 0.00 0.00 0.00 1.00 0.00 0.00 10 STEP_A -12.3 -1.2 -4.1 -2.0 -0.4 -0.2

2 bead O 1.000000 0.000000 0.000000 0.00 0.01 0.00 0.00 1.00 0.00 10 STEP_A -10.1 -1.0 -3.8 -1.7 -0.2 -0.1

END_FRAME

FRAME 1 time 0.010000

1 bead C 0.010000 0.000000 0.000000 0.01 0.00 0.00 1.00 0.00 0.00 10 STEP_B -12.0 -1.1 -4.0 -1.9 -0.4 -0.2

2 bead O 1.010000 0.000100 0.000000 0.01 0.01 0.00 0.00 1.00 0.00 10 STEP_B  -9.9 -1.0 -3.7 -1.6 -0.2 -0.1

END_FRAME


INFO

Replay stores dynamic bead positions and CG orientation vectors.

Static particles are stored once and reconstructed during playback.

TOTAL_FRAMES 2

DT 0.010000

\end{lstlisting}


\subsection{Data Model}


\begin{lstlisting}[language=C++]

struct XYZFullHeader {

    int version = 1;

    std::string units;

    std::string mode;

    std::vector<std::string> energy_layers;

    int static_count  = 0;

    int dynamic_count = 0;

};


struct XYZFullStatic {

    int id;

    std::string type;

    std::string sym;

    double x, y, z;

    double q;

    int cg_id;

};


struct XYZFullDynamic {

    int id;

    std::string type;

    std::string sym;

    double x, y, z;

    double vx, vy, vz;

    double ox, oy, oz;           // orientation vector

    int cg_id;

    std::string step_tag;

    std::array<double, 6> energy;

};


struct XYZFullFrame {

    int frame_index;

    double time;

    std::vector<XYZFullDynamic> dynamic;

};


struct XYZFullData {

    XYZFullHeader header;

    std::vector<XYZFullStatic>  static_particles;

    std::vector<XYZFullFrame>   frames;

    std::string info_block;

    int    total_frames;

    double dt;

};

\end{lstlisting}


\subsection{Reader / Writer API}


\begin{lstlisting}[language=C++]

namespace vsepr::io {


class XYZFullReader {

public:

    static XYZFullData read(const std::string& path);

};


class XYZFullWriter {

public:

    static void write(const std::string& path, const XYZFullData& data);

};


// Streaming writer: use for simulations longer than you can hold in RAM.

class XYZFullStreamWriter {

public:

    explicit XYZFullStreamWriter(const std::string& path);

    void write_header(const XYZFullHeader& header);

    void write_static_block(const std::vector<XYZFullStatic>& statics);

    void begin_frames();

    void write_frame(const XYZFullFrame& frame);

    void write_footer(const std::string& info,

                      int total_frames, double dt);

};


} // namespace vsepr::io

\end{lstlisting}


Use \type{XYZFullStreamWriter} for long simulations.  Accumulating 900\,000

frames in memory before writing them out is not a strategy.


\subsection{Replay Flow}


\begin{lstlisting}[language=C++]

void replay_xyzfull(const XYZFullData& data) {

    RenderWorld world;

    world.load_static(data.static_particles);

    for (const auto& frame : data.frames) {

        world.update_dynamic(frame.dynamic);

        world.update_cg_orientation(frame.dynamic);

        render_low_poly(world);

        sleep_fs(data.dt);

    }

}

\end{lstlisting}


\subsection{Shell Commands}


\begin{lstlisting}[language={},frame=single]

export xyzfull last    out.xyzfull

export xyzfull history replay.xyzfull

replay xyzfull         replay.xyzfull

tui    xyzfull         replay.xyzfull

convert xyzf xyzfull   trajectory.xyzf replay.xyzfull


# Options

export xyzfull history replay.xyzfull --energy-layers 6

export xyzfull history replay.xyzfull --dynamic-only

export xyzfull history replay.xyzfull --include-step-proxy

export xyzfull history replay.xyzfull --dt 0.01

\end{lstlisting}


\subsection{Validation Rules}


\begin{enumerate}

  \item File must begin with \texttt{XYZFULL}.

  \item \texttt{version} must be a supported value.

  \item Energy layer count must be 4, 5, or 6.

  \item Final two lines must be \texttt{TOTAL\_FRAMES \textit{N}} and

        \texttt{DT \textit{dt}}.

  \item Parsed frame count must equal \texttt{TOTAL\_FRAMES}.

  \item Dynamic IDs must be unique within each frame.

  \item Static IDs must not collide with dynamic IDs unless explicitly

        marked reusable.

  \item Orientation vector must satisfy $\|\mathbf{o}\| > 10^{-12}$;

        violation is warned, not fatal.

  \item Unknown \texttt{step\_tag} is permitted with a warning.

\end{enumerate}


\subsection{Implementation Files}


\begin{lstlisting}[language={},frame=single]

include/io/xyzfull_format.hpp

src/io/xyzfull_format.cpp

include/render/xyzfull_replay.hpp

src/render/xyzfull_replay.cpp

tests/test_xyzfull_read.cpp

tests/test_xyzfull_write.cpp

tests/test_xyzfull_footer.cpp

tests/test_xyzfull_replay.cpp

\end{lstlisting}


% ================================================

\section{Optional \texttt{INFO} Block and Compute Kernel Tags}

\label{sec:info_kernels}

% ================================================


Any format in the XYZ family may carry an optional \texttt{INFO} tail block.

In \texttt{.xyzfull} this block is structural (the file is invalid without

\texttt{TOTAL\_FRAMES} and \texttt{DT}).  For \texttt{.xyz}, \texttt{.xyza},

\texttt{.xyzc}, and \texttt{.xyzf}, the \texttt{INFO} block is \emph{optional}:

readers must not fail if it is absent.


\subsection{INFO Block Grammar (All Formats)}


\begin{lstlisting}[language={},frame=single]

INFO

<free text lines>

KERNEL  <kernel_id>  <backend>  <threads>  [<flags>]

KERNEL  ...

TOTAL_FRAMES <N>   # .xyzf and .xyzfull only

DT <dt>            # .xyzfull only; optional annotation elsewhere

END_INFO           # optional terminator; assumed at EOF

\end{lstlisting}


The \texttt{INFO} block begins at the token \texttt{INFO} on its own line and

extends to \texttt{END\_INFO}, or to end-of-file.  Free-text lines before the

first \texttt{KERNEL} directive are treated as human-readable notes.


\subsection{KERNEL Directive}


\begin{table}[htbp]

  \centering

  \caption{\texttt{KERNEL} directive fields.}

  \label{tab:kernel_fields}

  \begin{tabular}{lll}

    \toprule

    Field & Type & Description \\

    \midrule

    \texttt{kernel\_id} & string  & Logical kernel name (e.g.\ \texttt{nbody\_vdw}) \\

    \texttt{backend}    & string  & Execution backend: \texttt{cpu}, \texttt{cuda},

                                   \texttt{opencl}, \texttt{simd}, \texttt{tui} \\

    \texttt{threads}    & integer & Thread / warp count at time of write \\

    \texttt{flags}      & string* & Space-separated optional flags (see below) \\

    \bottomrule

  \end{tabular}

\end{table}


\subsubsection{Standard Kernel IDs}


\begin{table}[htbp]

  \centering

  \caption{Standard kernel identifiers.}

  \label{tab:kernel_ids}

  \begin{tabular}{lp{9cm}}

    \toprule

    Kernel ID & Description \\

    \midrule

    \texttt{nbody\_vdw}      & Pairwise van der Waals (Lennard-Jones or similar) \\

    \texttt{nbody\_coul}     & Pairwise Coulomb / electrostatics \\

    \texttt{bond\_harmonic}  & Harmonic bond stretching \\

    \texttt{angle\_cosine}   & Cosine angle bending \\

    \texttt{torsion\_fourier}& Fourier torsion series \\

    \texttt{thermostat\_nhc} & Nos\'{e}--Hoover chain thermostat \\

    \texttt{integrator\_vv}  & Velocity-Verlet integrator step \\

    \texttt{cg\_orient}      & CG orientation propagation (bead replay) \\

    \texttt{render\_proxy}   & Low-poly proxy reconstruction (\texttt{.xyzfull} only) \\

    \bottomrule

  \end{tabular}

\end{table}


\subsubsection{Optional Flags}


\begin{itemize}

  \item \texttt{pbc} --- periodic boundary conditions active during this kernel run

  \item \texttt{cutoff=<r>} --- non-bonded cutoff radius in \AA{}

  \item \texttt{skin=<r>} --- neighbour-list skin distance in \AA{}

  \item \texttt{dt=<val>} --- timestep used by this kernel (fs)

  \item \texttt{reduced} --- energy was written in reduced 4-layer mode

  \item \texttt{stream} --- file was produced by \type{XYZFullStreamWriter}

\end{itemize}


\subsection{INFO Block per Format}


\begin{table}[htbp]

  \centering

  \caption{\texttt{INFO} block applicability across formats.}

  \label{tab:info_formats}

  \begin{tabular}{lccccc}

    \toprule

    Field & \texttt{.xyz} & \texttt{.xyza} & \texttt{.xyzc} & \texttt{.xyzf} & \texttt{.xyzfull} \\

    \midrule

    Free-text notes       & opt. & opt. & opt. & opt. & opt. \\

    \texttt{KERNEL} lines & opt. & opt. & opt. & opt. & opt. \\

    \texttt{TOTAL\_FRAMES}& ---  & ---  & ---  & opt. & \textbf{req.} \\

    \texttt{DT}           & ---  & ---  & ---  & opt. & \textbf{req.} \\

    \texttt{END\_INFO}    & opt. & opt. & opt. & opt. & opt. \\

    \bottomrule

  \end{tabular}

\end{table}


For \texttt{.xyzc} the \texttt{KERNEL} block naturally records which kernels

were hot at checkpoint time, which aids reproducible continuation runs.

For \texttt{.xyzf} trajectory files, one \texttt{KERNEL} entry per integrator

pass is sufficient; do not emit one per frame.


\subsection{Example: \texttt{.xyzf} with INFO and Kernel Tags}


\begin{lstlisting}[language={},frame=single]

3

Step 0 | E = -12.30 kcal/mol | FIRE iter

C 0.000000 0.000000 0.000000

O 1.230000 0.000000 0.000000

H 2.100000 0.800000 0.000000

3

Step 1 | E = -13.11 kcal/mol | FIRE iter

C 0.001000 0.000000 0.000000

O 1.229000 0.000000 0.000000

H 2.099000 0.799000 0.000000


INFO

FIRE minimisation trace, 2 steps.

KERNEL  bond_harmonic   cpu   4   dt=0.010000

KERNEL  nbody_vdw       cpu   4   cutoff=12.0 skin=2.0

KERNEL  integrator_vv   cpu   4   dt=0.010000

TOTAL_FRAMES 2

END_INFO

\end{lstlisting}


\subsection{Example: \texttt{.xyzfull} INFO with Kernel Tags}


\begin{lstlisting}[language={},frame=single]

INFO

Replay stores dynamic bead positions and CG orientation vectors.

Static particles stored once; reconstructed during playback.

KERNEL  cg_orient       cpu    4   dt=0.010000

KERNEL  render_proxy    tui    1   stream

KERNEL  nbody_vdw       cuda  256  cutoff=12.0 pbc

TOTAL_FRAMES 2

DT 0.010000

\end{lstlisting}


\subsection{Reader Behaviour}


\begin{enumerate}

  \item If no \texttt{INFO} token is found: silently succeed; \texttt{KERNEL}

        list is empty.

  \item Unknown \texttt{KERNEL} IDs: stored as-is, warned in verbose mode.

  \item Unrecognised flags: stored as raw strings; no parse error.

  \item For \texttt{.xyzfull}: missing \texttt{TOTAL\_FRAMES} or \texttt{DT}

        is a hard error regardless of \texttt{INFO} presence.

\end{enumerate}


% ================================================

\section{Cross-References}

% ================================================


\begin{table}[htbp]

  \centering

  \caption{Related documentation.}

  \label{tab:xrefs}

  \begin{tabular}{lp{8cm}}

    \toprule

    Document & Relevance \\

    \midrule

    \file{docs/FILE\_FORMATS.md}              & Prose description of all I/O formats \\

    \file{section0\_identity\_state.pdf}       & Formal particle/cell/world definitions \\

    \file{section5\_integration.tex}           & MD timestep and integrator details \\

    \file{ALPHA\_MODEL\_BOOKLET.tex}           & Alpha model (Drude charge coupling) \\

    \file{section\_fuzzy\_ball\_model.tex}     & Coarse-grained bead visualisation \\

    \file{include/io/xyzfull\_format.hpp}      & \texttt{.xyzfull} reader/writer/stream API \\

    \file{include/render/xyzfull\_replay.hpp}  & Low-poly CG replay renderer \\

    \bottomrule

  \end{tabular}

\end{table}


\vfill


\begin{center}

\rule{0.5\textwidth}{0.4pt}\\[6pt]

\small Formation Engine Methodology v4.1.0 --- XYZ File Formats (Edition 3)

\end{center}


\end{document}
